%% kapa: exam=1 round=28 year=2017 period=1 form=A subject=civil_law subject_ko=민법 seq=11 number=11
%% source: Q-Net 감정평가사 기출문제 (cst00309) artlSeq=5199603 file=extracted/1008/5199603/2017년 제28회 감정평가사 1차시험 문제지/1교시 A형.pdf
%% conversion: kordoc 4.24.51 -> tools/gen1.py (auto); review=proofread-sonnet; flags=none
\begin{question}{11}
의사표시의 효력발생에 관한 설명으로 옳지 않은 것은? (다툼이 있으면 판례에 따름)
\begin{choices}
  \item 상대방 있는 의사표시는 상대방에게 도달한 때에 효력이 발생하는 것이 원칙이다.
  \item 의사표시의 상대방이 의사표시를 받은 때에 제한능력자인 경우에는 그 상대방의 법정대리인이 의사표시가 도달한 사실을 안 후라도 표의자는 그 의사표시로써 대항할 수 없다.
  \item 표의자가 의사표시의 통지를 발송한 후에 사망한 경우, 그 의사표시의 효력에 영향을 미치지 않는다.
  \item 내용증명우편이나 등기로 발송된 우편물은 반송 등의 특별한 사정이 없는 한 그 무렵 수취인에게 배달된 것으로 본다.
  \item 표의자가 과실 없이 상대방을 알지 못하거나 상대방의 소재를 알지 못하는 경우, 의사표시는 민사소송법의 공시송달의 규정에 의하여 송달할 수 있다.
\end{choices}
\end{question}
\answer{2}
